\chapter{外测度}
\label{chap:outer-measures}

测度只能在已经选定的 $\sigma$-代数上运算。实际构造通常反过来进行：先用
简单集合覆盖任意集合，取覆盖代价的下确界；再从全部子集中筛出与这个代价
相容的集合。外测度负责第一步，Carathéodory 条件负责第二步。

\section{覆盖与外测度}

\begin{definition}
映射 $\mu^*:\mathcal P(X)\to[0,+\infty]$ 称为 $X$ 上的
\term{外测度}，如果
\begin{enumerate}[label=\textup{(\roman*)}]
 \item $\mu^*(\varnothing)=0$；
 \item $A\subseteq B$ 时 $\mu^*(A)\le\mu^*(B)$；
 \item 对任意集合列，
 \[
  \mu^*\!\left(\bigcup_nA_n\right)\le\sum_n\mu^*(A_n).
 \]
\end{enumerate}
\end{definition}

外测度定义在全部子集上，却一般不具有可加性。它给出的是从外部覆盖所得
的代价；两个集合即使不交，也可能被同一简单集合同时覆盖，从而使并集的
外测度严格小于两者外测度之和。

最常用的构造如下。设 $\mathcal C\subseteq\mathcal P(X)$ 含空集，给每个
$C\in\mathcal C$ 指定代价 $\ell(C)\in[0,+\infty]$，且
$\ell(\varnothing)=0$。定义
\begin{equation}
 \mu^*(A)=\inf\left\{\sum_{n=1}^{\infty}\ell(C_n):
 A\subseteq\bigcup_{n=1}^{\infty}C_n, C_n\in\mathcal C\right\}.
 \label{eq:covering-outer-measure}
\end{equation}
若不存在覆盖，约定下确界为 $+\infty$。

\begin{proposition}
公式 \eqref{eq:covering-outer-measure} 定义外测度。
\end{proposition}

\begin{proof}
空集可由全为空集的集合列覆盖，故外测度为零。包含关系只扩大允许覆盖的
要求，给出单调性。对 $A\subseteq\bigcup_kA_k$，若右侧有限，给定
$\varepsilon>0$，为每个 $A_k$ 选覆盖，使其总代价小于
$\mu^*(A_k)+\varepsilon2^{-k}$。把所有覆盖集合排成一个序列便覆盖 $A$，
故
\[
 \mu^*(A)\le\sum_k\mu^*(A_k)+\varepsilon.
\]
令 $\varepsilon\downarrow0$。右侧无穷时不等式自动成立。
\end{proof}

\section{Carathéodory 可测性}

一个集合 $E$ 应当把任意测试集合 $A$ 的外测度准确分成落在 $E$ 内外的
两部分。次可加性已经给出“$\le$”方向，定义只要求反向不等式也成立。

\begin{definition}
设 $\mu^*$ 为外测度。集合 $E\subseteq X$ 称为
\term{Carathéodory 可测}，如果对每个 $A\subseteq X$ 都有
\begin{equation}
 \mu^*(A)=\mu^*(A\cap E)+\mu^*(A\setminus E).
 \label{eq:caratheodory-criterion}
\end{equation}
全部这类集合记为 $\mathcal M_{\mu^*}$。
\end{definition}

定义要对所有 $A$ 检验，是为了保证 $E$ 在任何后续集合运算中都能作为
不损失外测度的切割面。只检验 $A=X$ 不足以推出可数可加性。

\begin{theorem}[Carathéodory]
\label{thm:caratheodory-measurable-sets}
$\mathcal M_{\mu^*}$ 是 $\sigma$-代数，且
$\mu=\mu^*|_{\mathcal M_{\mu^*}}$ 是完备测度。
\end{theorem}

\begin{proof}
空集可测；条件 \eqref{eq:caratheodory-criterion} 对 $E$ 与 $E^c$ 完全
对称。若 $E,F$ 可测，对任意 $A$ 先按 $E$ 分割，再把 $A\cap E$ 按 $F$
分割，得到
\[
 \mu^*(A)\ge
 \mu^*(A\cap E\cap F)+\mu^*(A\cap E\setminus F)
 +\mu^*(A\setminus E).
\]
后两项的并是 $A\setminus(E\cap F)$；次可加性的反向整理给出
$E\cap F$ 可测。因此可测集构成代数。

设 $E_n$ 两两不交且均可测，记 $F_N=\bigcup_{n\le N}E_n$。反复分割给出
\[
 \mu^*(A)\ge
 \sum_{n=1}^N\mu^*(A\cap E_n)+\mu^*(A\setminus F_N).
\]
令 $E=\bigcup_nE_n$。由于 $A\setminus F_N\supseteq A\setminus E$，
再令 $N\to\infty$，有
\[
 \mu^*(A)\ge
 \sum_n\mu^*(A\cap E_n)+\mu^*(A\setminus E)
 \ge\mu^*(A\cap E)+\mu^*(A\setminus E).
\]
另一方向来自次可加性，所以 $E$ 可测。一般可数并先化为不交并，故
$\mathcal M_{\mu^*}$ 是 $\sigma$-代数。

在不交可测集列上，取 $A=\bigcup_nE_n$，上面的有限分割给出
$\mu^*(A)\ge\sum_{n\le N}\mu^*(E_n)$；令 $N\to\infty$，再与次可加性
合并，得到可数可加性。

若 $Z\subseteq N$ 且 $\mu^*(N)=0$，则对任意 $A$，
$\mu^*(A\cap Z)=0$，而 $A\setminus Z\subseteq A$。结合次可加性即得
$Z$ 满足 Carathéodory 条件，且测度为零。因此所得测度完备。
\end{proof}

\section{从预测度到外测度}

\begin{definition}
设 $\mathcal A_0$ 是 $X$ 上的代数。函数
$\mu_0:\mathcal A_0\to[0,+\infty]$ 称为\term{预测度}，如果
$\mu_0(\varnothing)=0$，并且每当两两不交的 $A_n\in\mathcal A_0$ 的并
$A=\bigcup_nA_n$ 仍属于 $\mathcal A_0$ 时，
\[
 \mu_0(A)=\sum_n\mu_0(A_n).
\]
\end{definition}

用 $\mathcal A_0$ 中集合覆盖任意 $E\subseteq X$，定义
\begin{equation}
 \mu^*(E)=\inf\left\{\sum_n\mu_0(A_n):
 E\subseteq\bigcup_nA_n, A_n\in\mathcal A_0\right\}.
 \label{eq:premeasure-induced-outer}
\end{equation}

\begin{lemma}
\label{lem:premeasure-agrees-on-algebra}
对 $A\in\mathcal A_0$，有 $\mu^*(A)=\mu_0(A)$，并且 $A$ 满足
Carathéodory 条件。
\end{lemma}

\begin{proof}
用 $A$ 自身覆盖给出 $\mu^*(A)\le\mu_0(A)$。反过来，若
$A\subseteq\bigcup_nA_n$，把 $A$ 依次分成
\[
 B_n=A\cap A_n\setminus\bigcup_{k<n}A_k.
\]
$B_n$ 两两不交、属于代数且并为 $A$。预测度的可数可加性和单调性给出
$\mu_0(A)=\sum_n\mu_0(B_n)\le\sum_n\mu_0(A_n)$。对所有覆盖取下确界。

固定 $E\subseteq X$ 和 $A\in\mathcal A_0$。任一代数覆盖 $(A_n)$ 可按
$A$ 分为 $A_n\cap A$ 与 $A_n\setminus A$，且
\[
 \mu_0(A_n)=\mu_0(A_n\cap A)+\mu_0(A_n\setminus A).
\]
对覆盖代价取下确界，得到
$\mu^*(E)\ge\mu^*(E\cap A)+\mu^*(E\setminus A)$，所以 $A$ 可测。
\end{proof}

\begin{theorem}[Carathéodory 延拓定理]
\label{thm:caratheodory-extension}
预测度 $\mu_0$ 在 $\sigma(\mathcal A_0)$ 上至少有一个测度延拓。由
\eqref{eq:premeasure-induced-outer} 构造的 Carathéodory 测度给出这样的
延拓，并且在整个 $\mathcal M_{\mu^*}$ 上已经完备。
\end{theorem}

\begin{proof}
公式 \eqref{eq:premeasure-induced-outer} 定义外测度。由引理
\ref{lem:premeasure-agrees-on-algebra}，
$\mathcal A_0\subseteq\mathcal M_{\mu^*}$，故
$\sigma(\mathcal A_0)\subseteq\mathcal M_{\mu^*}$。定理
\ref{thm:caratheodory-measurable-sets} 说明限制后的函数是测度；同一引理
说明它在 $\mathcal A_0$ 上等于 $\mu_0$。
\end{proof}

\section{延拓的唯一性}

存在性不自动给出唯一性。最简单的反例是只在代数
$\{\varnothing,X\}$ 上规定 $\mu_0(X)=+\infty$；任何总质量无穷的测度
都延拓它。$\sigma$-有限性排除了这种完全没有局部有限控制的情形。

唯一性证明要用到一个集合族闭包结论。称包含 $X$、对补集以及两两不交的
可数并封闭的集合族为 \term{Dynkin 系统}；称对有限交封闭的集合族为
\term{$\pi$-系统}。

\begin{lemma}[$\pi$--$\lambda$ 定理]
\label{lem:pi-lambda}
若 $\mathcal P$ 是 $X$ 上的 $\pi$-系统，$\mathcal D$ 是包含
$\mathcal P$ 的 Dynkin 系统，则
\[
 \sigma(\mathcal P)\subseteq\mathcal D.
\]
\end{lemma}

\begin{proof}
记 $\lambda(\mathcal P)$ 为包含 $\mathcal P$ 的最小 Dynkin 系统。先
固定 $A\in\mathcal P$，令
\[
 \mathcal D_A=\{B\subseteq X:A\cap B\in\lambda(\mathcal P)\}.
\]
因为 $A\in\lambda(\mathcal P)$，集合族 $\mathcal D_A$ 包含 $X$；
Dynkin 系统的差集封闭性和不交可数并封闭性说明 $\mathcal D_A$ 仍是
Dynkin 系统。$\mathcal P$ 对有限交封闭，所以
$\mathcal P\subseteq\mathcal D_A$，从而
$\lambda(\mathcal P)\subseteq\mathcal D_A$。

再固定 $B\in\lambda(\mathcal P)$，令
\[
 \mathcal E_B=\{A\subseteq X:A\cap B\in\lambda(\mathcal P)\}.
\]
前一段说明 $\mathcal P\subseteq\mathcal E_B$；同样的检验表明
$\mathcal E_B$ 是 Dynkin 系统。因此
$\lambda(\mathcal P)\subseteq\mathcal E_B$，即
$\lambda(\mathcal P)$ 对有限交封闭。一个对有限交封闭的 Dynkin 系统
同时对补集和任意可数并封闭，因而是 $\sigma$-代数。于是
\[
 \lambda(\mathcal P)=\sigma(\mathcal P).
\]
又 $\lambda(\mathcal P)\subseteq\mathcal D$，结论得证。
\end{proof}

\begin{theorem}[延拓唯一性]
\label{thm:premeasure-extension-uniqueness}
若 $X=\bigcup_nX_n$，其中 $X_n\in\mathcal A_0$ 且
$\mu_0(X_n)<\infty$，则 $\mu_0$ 在 $\sigma(\mathcal A_0)$ 上的测度延拓
唯一。
\end{theorem}

\begin{proof}
先设 $\mu_0(X)<\infty$，且 $\mu,\nu$ 是两个延拓。令
\[
 \mathcal D=\{E\in\sigma(\mathcal A_0):\mu(E)=\nu(E)\}.
\]
由于两测度在 $X$ 上有限且相等，$E\in\mathcal D$ 推出
$E^c\in\mathcal D$；对两两不交的 $E_n\in\mathcal D$，可数可加性给出
$\bigcup_nE_n\in\mathcal D$。所以 $\mathcal D$ 是 Dynkin 系统。代数
$\mathcal A_0$ 对有限交封闭，并包含于 $\mathcal D$。引理
\ref{lem:pi-lambda} 说明
$\sigma(\mathcal A_0)\subseteq\mathcal D$。

一般情形把 $X_n$ 替换为两两不交的有限测度代数成员。对每个 $n$，受限
测度 $E\mapsto\mu(E\cap X_n)$ 与 $E\mapsto\nu(E\cap X_n)$ 在
$\mathcal A_0$ 上相等且有限，由前一段在生成 $\sigma$-代数上相等。最后
\[
 \mu(E)=\sum_n\mu(E\cap X_n)
 =\sum_n\nu(E\cap X_n)=\nu(E).
\]
\end{proof}

$\pi$--$\lambda$ 定理以后还会用于乘积测度和概率分布的唯一性。习题要求
读者不参照正文重新写出闭包论证。

\section*{习题}
\addcontentsline{toc}{section}{习题}

\subsection*{A组　外测度与可测性}
\begin{enumerate}[label=\textbf{\thechapter.\arabic*.}]
 \exerciseitem{outer-measure-covering-proof}
 完整证明覆盖公式 \eqref{eq:covering-outer-measure} 的可数次可加性。
 \exerciseitem{outer-null-measurable}
 证明每个外测度零集及其任意子集都满足 Carathéodory 条件。
 \exerciseitem{caratheodory-finite-unions}
 直接验证 Carathéodory 可测集对有限并和差封闭。
 \exerciseitem{caratheodory-countable-additivity}
 补全定理 \ref{thm:caratheodory-measurable-sets} 中可数可加性的有限分割
 极限步骤。
 \exerciseitem{metric-outer-measure-borel}
 设 $X$ 为度量空间，外测度 $\mu^*$ 对距离为正的集合满足
 $\mu^*(A\cup B)=\mu^*(A)+\mu^*(B)$。证明所有闭集均 Carathéodory 可测，
 因而 $\mathcal B(X)\subseteq\mathcal M_{\mu^*}$。
\end{enumerate}

\subsection*{B组　预测度与延拓}
\begin{enumerate}[label=\textbf{\thechapter.\arabic*.},resume]
 \exerciseitem{premeasure-outer-agreement}
 逐项补全引理 \ref{lem:premeasure-agrees-on-algebra}。
 \exerciseitem{finite-partition-premeasure}
 给定有限分割 $X=A_1\dot\cup\cdots\dot\cup A_m$ 和权重 $w_j\ge0$，在
 由这些原子生成的代数上构造预测度，并求其 Carathéodory 延拓。
 \exerciseitem{extension-nonunique-infinite-total}
 在 $X=\N$ 上，从代数 $\{\varnothing,X\}$ 和规定
 $\mu_0(X)=+\infty$ 出发，写出两个不同的测度延拓。
 \exerciseitem{sigma-finite-disjoint-cover}
 证明延拓唯一性定理中的 $X_n$ 可替换为两两不交的代数成员，并保持有限
 测度。
 \exerciseitem{pi-lambda-theorem}
 证明：若 $\mathcal P$ 是对有限交封闭的集合族，$\mathcal D$ 是包含
 $\mathcal P$ 的 Dynkin 系统，则 $\sigma(\mathcal P)\subseteq\mathcal D$。
\end{enumerate}

\subsection*{C组　构造方法}
\begin{enumerate}[label=\textbf{\thechapter.\arabic*.},resume]
 \exerciseitem{outer-measure-from-existing-measure}
 设 $(X,\mathcal A,\mu)$ 为 $\sigma$-有限测度空间，并定义
 $\mu^*(A)=\inf\{\mu(E):A\subseteq E, E\in\mathcal A\}$。证明它是外测度，
 并证明其 Carathéodory $\sigma$-代数恰为原测度空间的完备化。指出证明中
 $\sigma$-有限性用在何处。
 \exerciseitem{outer-measure-restriction}
 若 $Y\subseteq X$，证明 $\nu^*(A)=\mu^*(A)$（$A\subseteq Y$）是 $Y$
 上的外测度，并比较两者的可测集。
 \exerciseitem{premeasure-finite-cover-subadditivity}
 设 $A,A_1,\ldots,A_m\in\mathcal A_0$ 且
 $A\subseteq\bigcup_{j=1}^mA_j$。证明
 $\mu_0(A)\le\sum_j\mu_0(A_j)$。
 \optionalexerciseitem{caratheodory-maximal-domain}
 证明：若一个测度 $\nu$ 定义在某个 $\sigma$-代数 $\mathcal C$ 上，且
 $\nu^*$ 是由 $\nu$ 的可测覆盖得到的外测度，则
 $\mathcal C\subseteq\mathcal M_{\nu^*}$；说明完备化为何自然包含其中。
\end{enumerate}
